\documentclass{amsart}
\usepackage{amsmath, amssymb}
\title{Math 21b --- Worksheet 2: Eigenvalues, Eigenvectors, and Diagonalization}
\date{}
\begin{document}
\maketitle

\textbf{Warm-up.} Let $A = \begin{pmatrix}3&1\\0&2\end{pmatrix}$. Verify that $v = \begin{pmatrix}1\\0\end{pmatrix}$ is an eigenvector. What is the corresponding eigenvalue?

\bigskip

\textbf{Exercise 1.} For each matrix, find all eigenvalues and a basis for each eigenspace.
\[
(a)\; \begin{pmatrix}4&-2\\1&1\end{pmatrix}
\qquad
(b)\; \begin{pmatrix}2&1&0\\0&2&0\\0&0&3\end{pmatrix}
\]

\bigskip

\textbf{Exercise 2.} Determine whether each matrix in Exercise 1 is diagonalizable. If so, write $A = PDP^{-1}$ explicitly.

\bigskip

\textbf{Exercise 3.} Use diagonalization to compute $A^{10}$ for
\[
A = \begin{pmatrix}1&2\\0&3\end{pmatrix}.
\]

\bigskip

\textbf{Exercise 4 (Application --- Population Dynamics).} A population of rabbits and foxes evolves each year according to $x_{t+1} = Ax_t$ where
\[
A = \begin{pmatrix}1.2 & -0.1\\0.05 & 0.9\end{pmatrix}.
\]
Starting from $x_0 = \begin{pmatrix}100\\20\end{pmatrix}$:
\begin{enumerate}
    \item[(a)] Find the eigenvalues. Do both populations grow, shrink, or have mixed behavior?
    \item[(b)] Find the long-run trajectory.
\end{enumerate}

\bigskip

\textbf{Reflection.} Why might a matrix have complex eigenvalues even though all its entries are real? What does this mean geometrically?

\end{document}
